%% ==============================
\chapter{Einleitung}
\label{cha:einleitung}
%% ==============================

\section{Motivation}
\label{sec:einleitung-motivation}

Digitale Technologien haben in den letzten Jahren sowohl die Effizient und Angemessenheit der medizinischen Versorgung in Krankenhäusern als auch die Wirksamkeit einzelner Interventionen erhöht haben~\cite{ibrahim2022, keasberry2017}. Dennoch bleibt die Verwaltung der dabei entstehenden Gesundheitsdaten seit Jahren eine der zentralen Herausforderungen der Digitalisierung im Gesundheitswesen~\cite{stachwitz2023, bratan2022}. Krankenhäuser, medizinische Fachkräfte, Patienten und weitere Akteure verwenden heterogene IT-Systeme, die sich in Standards, Technologien und medizinischen Terminologien unterscheiden; dies erschwert den konsistenten Datenaustausch zwischen Institutionen~\cite{adel2022, torabmiandoab2023}. Gleichzeitig verschärft die Verwendung elektronischer Patientenakten die Sicherheitsanforderungen, da der Zugriff von vielen unterschiedlichen Stakeholdern erfolgt~\cite{mehrotra2012}. Gesundheitsdaten gelten gemäß der Datenschutz-Grundverordnung als eine besondere Kategorie personenbezogener Daten und erfordern daher einen hohen Schutz: Es muss sichergestellt sein, dass nur autorisierte Personen Zugriff erhalten und dass jeder Zugriff eindeutig kontrollierbar bleibt. Hinzu kommen hohe Anforderungen an Performanz, Verfügbarkeit und Ausfallsicherheit, da technische Störungen direkt den Zugriff auf Informationen, die für die Behandlung relevant sind, beeinträchtigen. Die Einführung der elektronischen Patientenakte (ePA) in Deutschland~\cite{bmg2026} zeigt beide Seiten gleichzeitig: Mehr als hundert Millionen Dokumentewurden innerhalb eines Jahres hinterlegt~\cite{gematik2026}, während bereits der verpflichtende Start von technischen Störungen begleitet wurde~\cite{apothekeadhoc2025}.

Zur Verbesserung eines solchen dezentralen Datenmanagements werden Distributed-Ledger-Technologien (DLT) als möglicher architektonischer Baustein diskutiert~\cite{ferreira2024}. Ihre Relevanz liegt weniger in der Dezentralisierung als solcher, sondern in der Bereitstellung einer gemeinsamen Vertrauensschicht zwischen unabhängigen Akteuren, ohne dass eine zentrale Instanz als einziger Kontrollpunkt verwendet werden muss~\cite{ferdous2020}. Auf dieser Grundlage kann DLT zur Verwaltung von Identitäten, Zugriffsrechten und Einwilligungen eingesetzt werden und kann die Architektur gegenüber dem Ausfall einzelner Komponenten robuster machen. Die bestehenden Arbeiten basieren hauptsächlich auf blockchainbasierte Ansätze, die Nachvollziehbarkeit, Manipulationsresistenz und eine dezentrale Abstimmung zwischen mehreren Beteiligten versprechen~\cite{sonkamble2021, reegu2023}.

Nichtsdestotrotz zeigen sich in der Literatur wiederkehrende Einschränkungen solcher Lösungen: Skalierbarkeitsprobleme~\cite{mazlan2020}, begrenzte Performanz~\cite{singh2021}, zusätzliche Kosten für Speicherung und Wartung~\cite{kasyapa2024} sowie ein erhöhter Energieaufwand der Konsensfindung, der besonders bei rechenintensiven Verfahren wie Proof-of-Work signifikant ist, sich aber nicht auf alle blockchainbasierten Architekturen gleichermaßen übertragen lässt~\cite{bada2021}. Außerdem gibt es eine domänenspezifische Einschränkung: Sensible Gesundheitsdaten sind aufgrund der Datenschutzanforderungen und der eingeschränkten Löschbarkeit für eine On-Chain-Speicherung ungeeignet, und je mehr Daten auf dem Ledger abgelegt werden, desto stärker werden Performanz und Speicherbedarf beeinflusst. Daraus ergibt sich ein Zielkonflikt zwischen manipulationsresistenter dezentraler Speicherung und einer datenschutzkonformen, effizienten Architektur - und damit ein Anlass, Alternativen zu prüfen, die über die Blockchain hinausgehen.

Eine solche Alternative sind hashgraphbasierte DLT-Ansätze. Sie speichern Transaktionen nicht als lineare Blockkette, sondern als Ereignisse in einem gerichteten azyklischen Graphen und erreichen Konsens über die Mechanismen Gossip about Gossip und Virtual Voting, bei denen die Knoten Kommunikationsmetadaten verbreiten und die Abstimmung daraus lokal rekonstruieren, statt sie explizit durchzuführen~\cite{baird2016}. Die Manipulationsresistenz ergibt sich dabei aus der kryptographischen Verkettung der Ereignisse: Jedes Ereignis verweist über Hashwerte auf vorangegangene Ereignisse und ist signiert, sodass nachträgliche Änderungen die Hashreferenzen und Signaturen ungültig machen und damit erkennbar werden. Da auf energie- und rechenintensive Mining-Prozesse verzichtet wird und die graphenbasierte Struktur eine parallele Verarbeitung der Transaktionen ermöglicht, werden hashgraphbasierte Ansätze in der Literatur mit Vorteilen hinsichtlich Durchsatz und Kosten in Verbindung gebracht~\cite{akhtar2019}.

Im Kontext des dezentralen Gesundheitsdatenmanagements sind solche Ansätze bisher jedoch deutlich seltener untersucht worden. Eine gezielte Suche über Google Scholar verdeutlicht dieses Ungleichgewicht: Die Anfrage \texttt{\textquotedbl blockchain\textquotedbl\ AND \textquotedbl healthcare\textquotedbl\ AND \textquotedbl data management\textquotedbl} liefert etwa 91.300 Treffer, die entsprechende Anfrage mit \texttt{"hashgraph"} dagegen nur etwa 492; bei \texttt{\textquotedbl electronic health records\textquotedbl} stehen rund 39.100 Treffer etwa 195 Treffern gegenüber. Diese Beobachtung ist lediglich ein erster Indikator; die systematische Literaturrecherche im Kapitel~\ref{cha:state_of_the_art} bestätigt und präzisiert sie: Im Gesundheitskontext existiert bislang keine Architektur, die ein SSI-basiertes Identitätsmodell auf einem hashgraphbasierten Ledger realisiert, und die publizierten Evaluationsergebnisse sind aufgrund unterschiedlicher Messgrößen, Lastprofile und Systemgrenzen untereinander nicht vergleichbar~\cite{sehimibendaoud2023, ferreira2024}. Ein belastbarer Vergleich beider Ledger-Familien lässt sich daher nicht aus der vorhandenen Literatur ableiten, sondern erfordert eine eigene Messung beider Varianten unter identischen Bedingungen.

Einen konkreten Ausgangspunkt für eine solche Untersuchung bietet die Referenzarchitektur von Erler et al.~\cite{erler2023}. Sie beschreibt eine sicherheitsgetrieben entworfene Anwendung für dezentrales Gesundheitsdatenmanagement, in der die Gesundheitsdaten selbst off-chain in den Systemen der erzeugenden Einrichtung verbleiben und der Ledger ausschließlich als öffentlicher Vertrauensanker für Identitäten, Schemas, Credential Definitions und Widerrufsinformationen dient~\cite{erler2023, erler2024}. Die Ledger-Schicht ist damit klar von der übrigen Architektur abgegrenzt, sodass die zugrunde liegende DLT ausgetauscht werden kann, ohne die Agenten- und Kommunikationsschicht zu verändern - und genau diese Eigenschaft ist die Voraussetzung dafür, gemessene Unterschiede tatsächlich der Ledger-Technologie zuzuordnen.

Dieser Austausch ist allerdings nicht nur ein Wechsel der Implementierung, sondern ein Wechsel der Systemklasse. Hyperledger Indy ist ein Zustands-Ledger: Er hält den aktuellen Zustand der registrierten Identitätsobjekte selbst vor und setzt Schreibberechtigungen über die Rollen Steward und Endorser unmittelbar im Ledger durch~\cite{torongo2023}. Der Hedera Consensus Service ist dagegen ein reiner Ordnungsdienst: Er ordnet Nachrichten innerhalb eines Topics manipulationsresistent und zeitlich eindeutig, interpretiert ihren Inhalt jedoch nicht und leitet daraus keinen Zustand ab. Alles, was Indy im Ledger selbst erbringt - Zustandshaltung, Objektsemantik und Berechtigungsprüfung - muss auf dem Hedera Consensus Service oberhalb des Ledgers rekonstruiert werden. Diese Verlagerung bestimmt sowohl die zu treffenden Entwurfsentscheidungen als auch die Unterschiede, die eine Evaluation überhaupt messbar machen soll, und bildet damit den technischen Kern dieser Arbeit.

\section{Zielsetzung und Forschungsfragen}
\label{sec:einleitung-zielsetzung}

Das Ziel dieser Arbeit ist es, eine hashgraphbasierte Alternative zur Ledger-Schicht der genannten Referenzarchitektur zu konzipieren, prototypisch umzusetzen und gegen die blockchainbasierte Ausgangsvariante zu evaluieren. Konkret wird die auf Hyperledger Indy basierende Verifiable Data Registry durch eine auf dem Hedera Consensus Service basierende Umsetzung ersetzt, während die Agenten- und Kommunikationsschicht unverändert bleibt. Beide Varianten werden anschließend in einem Proof-of-Concept unter identischen Bedingungen entlang der drei Dimensionen Performanz, Skalierbarkeit und Sicherheit gemessen und verglichen. Diese Dimensionen werden gewählt, weil hashgraphbasierten Ansätzen in der Literatur gerade hinsichtlich Transaktionsdurchsatz, Latenz und Rechenaufwand potenzielle Vorteile zugeschrieben werden~\cite{akhtar2019}, und weil sie zugleich diejenigen Eigenschaften abdecken, an denen sich die Eignung einer Ledger-Schicht für den Gesundheitskontext entscheidet.

Die Arbeit untersucht entsprechend die folgende Hauptforschungsfrage: Inwieweit eignen sich blockchain- und hashgraphbasierte Distributed-Ledger-Technologien als Verifiable Data Registry einer SSI-basierten Architektur für dezentrales Gesundheitsdatenmanagement hinsichtlich Performanz, Skalierbarkeit und Sicherheit? Daraus werden drei Teilfragen abgeleitet, die den drei Evaluationsdimensionen entsprechen:

\begin{description}
  \item[FF1 - Performanz:] Welche Unterschiede bestehen zwischen der Indy- und der HCS-basierten Ledger-Schicht hinsichtlich Durchsatz, operationsspezifischer Latenz und Ressourcenverbrauch?
  \item[FF2 - Skalierbarkeit:] Wie verhalten sich beide Varianten bei steigender Nebenläufigkeit sowie bei wachsendem Bestand verwalteter Identitätsobjekte?
  \item[FF3 - Sicherheit:] Welche Unterschiede weisen beide Varianten hinsichtlich Konsens-Fehlertoleranz, Durchsetzbarkeit der Schreibautorisierung und Bedrohungsresistenz auf?
\end{description}


\section{Methodisches Vorgehen und Aufbau der Arbeit}
\label{sec:einleitung-aufbau}

Methodisch gliedert sich die Arbeit in vier Schritte. Zuerst wird eine strukturierte Literaturrecherche durchgeführt, aus der die Forschungslücke sowie die Anforderungen, Architekturbausteine und Vergleichskriterien für die spätere Gegenüberstellung abgeleitet werden. Auf dieser Grundlage wird die Ledger-Schicht der Referenzarchitektur analysiert und ein Konzept für die hashgraphbasierte Alternative entworfen, das jede Entwurfsentscheidung explizit gegen die zuvor formulierten Anforderungen abwägt. Anschließend wird dieses Konzept als Proof-of-Concept implementiert, indem die Ledger-Schicht der bestehenden Referenzimplementierung ausgetauscht wird. Abschließend werden beide Varianten anhand der definierten Kriterien gemessen und die Ergebnisse eingeordnet.

Die Arbeit gliedert sich entsprechend wie folgt: Das Kapitel~\ref{cha:grundlagen} führt die konzeptionellen Grundlagen ein - die Blockchain-Technologie mit dem Hyperledger-Indy/Aries-Ökosystem, den Hashgraph-Konsens mit dem Hedera Consensus Service, die Konzepte der Self-Sovereign Identity sowie eine Gegenüberstellung der Speichermodelle beider Technologien. Das Kapitel~\ref{cha:state_of_the_art} stellt den Stand der Wissenschaft und Technik dar, ordnet die identifizierte Literatur taxonomisch ein, leitet daraus die Forschungslücke und die Vergleichskriterien ab und beschreibt abschließend die Referenzarchitektur im Detail. Das Kapitel~\ref{cha:konzept} analysiert die Ledger-Schicht dieser Architektur, leitet die funktionalen und nicht-funktionalen Anforderungen ab, trifft die Entwurfsentscheidungen für die HCS-basierte Alternative und legt das Evaluationskonzept fest. Das Kapitel~\ref{cha:implementierung} dokumentiert die prototypische Umsetzung und den Messaufbau, das Kapitel~\ref{sec:results} berichtet die Messergebnisse entlang der Lastszenarien, und das Kapitel~\ref{sec:discussion} diskutiert diese Ergebnisse und ordnet sie in die Gültigkeitsgrenzen der Evaluation ein. Das Kapitel~\ref{sec:conclusion} fasst die Arbeit zusammen, beantwortet die Forschungsfragen und gibt einen Ausblick auf weiterführende Untersuchungen.